\begin{textsample}{POS Dim 3 – vem\_ed\_16 – Score 6.00 – vem\_ed\_16/t071.txt}  \label{ex:f3_pos_018}
ACONTECE Programe-se para os eventos de 2023 relativos aos Intercâmbios Virtuais: FAUBAI 2023 Conference — 15 a 19 de abril A Associação Brasileira de Educação Internacional (FAUBAI) retoma o formato presencial, em Belo Horizonte, celebrando seus 35 anos de existência com o tema “Building knowledge with all voices”.
Submissões encerradas.
Programa e informações em: https://faubai.org.br/conf/2023/ IVES Cesu 2023 — 25 de maio O evento International Virtual Exchange Simposium, organizado pela equipe dos PCIs / Cesu, terá inscrições no Even3 e será transmitido pelo Canal Cesu no YouTube: https://www.youtube.com/@CanalCESU Red Latam COIL — 12 a 16 de junho O terceiro congresso da Red Latinoamericana COIL ocorrerá on-line.
Submissão de propostas até 12 de março: https://www.uv.mx/coil/conferencia-anual/2023/ CBTecle — 13 a 15 de setembro Realizado a cada dois anos pela Coordenação de Línguas e Projetos Internacionais da Cesu / CPS, o Congresso Brasileiro de Línguas na Formação Técnica e Tecnológica retoma o formato presencial, neste ano na Fatec Guaratinguetá: https://cesu.cps.sp.gov.br/cbt ecle/ Jornada de PCIs — 5 de \textbf{outubro} Evento organizado pela equipe dos PCIs / Cesu, com apresentação de relatos de experiências de PCIs nas Fatecs.
Inscrições pelo Even3 e transmissão pelo Canal Cesu no YouTube: https://www.youtube.com/@CanalCESU IVEC — de 30 de \textbf{outubro} a 01 de novembro A principal conferência mundial sobre Intercâmbios Virtuais ocorrerá pela primeira vez no Brasil, será sediada em São Paulo.
Submissão de resumos até 15 de março: https://iveconference.org/ PCIs em artigos na Revista CBTecLE É notável o crescimento da pesquisa com foco nos Projetos Colaborativos Internacionais.
Veja três artigos a respeito dessa \textbf{temática}, publicados na edição de dezembro/2022 da Revista CBTecLE, que \textbf{conquistou} a classificação B1 pelo Qualis-Capes: 1) El desarrollo de competencias interculturales en el proceso de enseñanza-aprendizaje de lenguas en modelo COIL: más que una experiencia un reto – Regiane Souza Camargo Moreira (Fatec Guaratinguetá) e Mónica Díaz (Uniminuto) 2) Internacionalización: aprendizaje intercultural en entornos virtuales – Elizabeth Colorado Herrera (Fatec Itaquaquecetuba) 3) Explorando a noção de qualidade em um intercâmbio virtual: um relato de experiência - Patrícia Sales Patrício (Fatec Ipiranga), Danilo Nunes (Fatec Praia Grande), Luciana Maria Gasparelo Spigolon (Fatec Ribeirão Preto) e Zulmira Rodrigo Torrecilhas (Fatec Pindamonhangaba). https://revista.cbtecle.com.br/

% matched lemmas: conquistar, outubro, temática
\end{textsample}
